\documentclass[12pt]{article}
\usepackage[T1]{fontenc}
\usepackage[utf8]{inputenc}
\usepackage{lmodern}
\usepackage{textcomp}
\usepackage{enumitem}
\usepackage{geometry}
\geometry{margin=1in}
\begin{document}
\begin{center}
Answer Key - Constitutional Law - Graduate Level Assessment 4 \
\small (Answers are ordered exactly as in the question paper.)
\end{center}

\section*{Multiple Choice}
\begin{enumerate}[label=\arabic*.]
\item \textbf{Answer:} B
\\ \textit{Options:} A) No duty for artificial conditions.; B) A duty for natural conditions.; C) No duty for natural conditions.; D) No duty for active operations.
\\ \textit{Note:} Correct option: B - A duty for natural conditions.
\item \textbf{Answer:} C
\\ \textit{Options:} A) There is no distinction between the two forms of legal reasoning.; B) Judges are appointed to interpret the law, not to make it.; C) It is misleading to pigeon-hole judges in this way.; D) Judicial reasoning is always formal.
\\ \textit{Note:} Correct option: C - It is misleading to pigeon-hole judges in this way.
\item \textbf{Answer:} A
\\ \textit{Options:} A) Savigny; B) Kelsen; C) Austin; D) Hart; E) Bentham
\\ \textit{Note:} Correct option: A - Savigny
\item \textbf{Answer:} A
\\ \textit{Options:} A) The Nazi law in question was validly enacted.; B) The Nazi rule of recognition was unclear.; C) Hart agreed with Fuller's argument.; D) The court misunderstood the legislation.; E) Fuller misconstrued the purpose of the law.
\\ \textit{Note:} Correct option: A - The Nazi law in question was validly enacted.
\item \textbf{Answer:} A
\\ \textit{Options:} A) Command, sovereign and legal remedy; B) Command, sovereign and authority; C) Command, sovereign and legal enforcement; D) Command, sovereign and legal responsibility; E) Command, sovereign and obedience by subject
\\ \textit{Note:} Correct option: A - Command, sovereign and legal remedy
\end{enumerate}

\section*{True / False}
\begin{enumerate}[label=\arabic*.]
\setcounter{enumi}{5}
\item \textbf{Answer:} False
\\ \textit{Options:} A) True; B) False
\\ \textit{Note:} LLM-generated false statement. Correct answer: 'Cold pursuit'.
\item \textbf{Answer:} True
\\ \textit{Options:} A) True; B) False
\\ \textit{Note:} LLM-generated true statement based on MCQ.
\item \textbf{Answer:} True
\\ \textit{Options:} A) True; B) False
\\ \textit{Note:} LLM-generated true statement based on MCQ.
\item \textbf{Answer:} True
\\ \textit{Options:} A) True; B) False
\\ \textit{Note:} LLM-generated true statement based on MCQ.
\item \textbf{Answer:} False
\\ \textit{Options:} A) True; B) False
\\ \textit{Note:} LLM-generated false statement. Correct answer: 'Prior notice and prior evidentiary hearing'.
\end{enumerate}

\section*{Long Form Response}
\begin{enumerate}[label=\arabic*.]
\setcounter{enumi}{10}
\item \textbf{Answer:} 50\% to Joan and the income from the remaining 50\% to Joan for life, remainder to the Salvation Army, if Joan files a timely petition protesting the devise to the Salvation Army.. Explain why this is the correct answer and provide supporting evidence. Reference key facts or concepts that support this choice.
\\ \textit{Note:} Long-form derived from MCQ; award credit for stating the correct idea and giving supporting reasoning.
\item \textbf{Answer:} A bidder may retract his bid at any time until the falling of the hammer.. Explain why this is the correct answer and provide supporting evidence. Reference key facts or concepts that support this choice.
\\ \textit{Note:} Long-form derived from MCQ; award credit for stating the correct idea and giving supporting reasoning.
\end{enumerate}

\vfill
\noindent\textit{End of Answer Key}
\end{document}